\documentclass[11pt]{article}
\usepackage{series}
\usepackage{amsmath,amssymb,amsthm}
\usepackage{geometry}
\usepackage{hyperref}
\usepackage{enumitem}

\geometry{margin=1in}


% ---------- metadata ----------
\title{The Modal Triplet Theory Program B2:\\ Gauge Structure as Redundancy Encoding\\
in the Modal Triplet Theory Program}
\author{Peter Nero}
\date{January 2026}

\begin{document}

\maketitle

\begin{abstract}
We show that gauge structure arises as a necessary encoding class once local
describability and overlap consistency coexist with lens obstructions, i.e.\ with
non-unique but consistent local representations. In the Modal Triplet Theory
(MTT) framework, gauge is not postulated as a symmetry of underlying reality and
is not introduced as a field-theoretic axiom. Rather, it is identified as
bookkeeping of descriptive redundancy: the unique encoding that organizes and
compensates for the non-uniqueness of admissible local sections while preserving
overlap consistency and kinematic persistence.

We define gauge transformations as admissible automorphisms of encoding fibers,
and show that gauge freedom is unavoidable whenever lens obstructions are
present and global canonical representatives do not exist. We further show that
in smooth realization classes, redundancy bookkeeping is realized by bundle
structure with a connection-like object implementing consistent comparison of
redundant representatives. This connection is distinct in role from the
kinematic consistency encoding identified with gravity: gravity resolves circle
obstructions (loop inconsistency), while gauge resolves lens obstructions
(non-unique lifts). The analysis clarifies the structural origin of gauge
universality and prepares the ground for quantization as the encoding response
to nil obstructions.
\end{abstract}

\tableofcontents
\newpage

% ============================================================
\section{Introduction and Scope}

The structural core of Modal Triplet Theory establishes that reduced
descriptions exist only locally on admissible domains, that such descriptions
must be related by controlled re-encoding on overlaps, and that no single global
reduced encoding exists. Coherent kinematics then defines motion and causal
structure as persistence across overlapping admissible encodings, without
assuming spacetime or dynamical laws. A subsequent classification identifies
three and only three obstruction types to global coherence: circle, lens, and
nil.

The purpose of the present paper is to analyze the structural consequences of
\emph{lens} obstructions. A lens obstruction is the failure of global coherence
by \emph{redundancy}: multiple admissible local representatives exist that are
mutually consistent on overlaps, but no canonical global choice exists. In such
regimes, kinematic persistence and overlap consistency remain intact, but the
reduced description is not unique. This non-uniqueness is neither contradiction
(circle) nor termination of description (nil). It is an independent obstruction
type and therefore demands its own encoding response.

We show that lens obstructions force the introduction of an additional encoding
class whose sole role is to organize descriptive redundancy. This encoding class
is what is conventionally identified as gauge structure. In MTT terms, gauge is
not a fundamental physical symmetry; it is the unique bookkeeping structure that
compensates for the non-uniqueness of admissible local sections while preserving
consistent reduced description.

Throughout this paper we adhere to three principles:

\begin{itemize}
\item Gauge is an \emph{encoding} of redundancy, not a postulated interaction or
field.
\item Gauge transformations are \emph{automorphisms of representation}, not
transformations of underlying reality.
\item Geometric objects commonly associated with gauge theory (bundles,
connections) appear only as \emph{realizations} of the redundancy encoding, not
as axioms of the structural theory.
\end{itemize}

We emphasize the separation between gauge and gravity in the MTT program.
Gravity, treated elsewhere, arises as kinematic consistency bookkeeping required
to resolve circle obstructions (loop-dependent inconsistency). Gauge arises here
as redundancy bookkeeping required to resolve lens obstructions (non-unique but
consistent representation). The two encodings are structurally distinct, though
they may couple in particular realizations.

\begin{remark}[Position in the series]
This paper depends only on the structural core, coherent kinematics, and the
circle--lens--nil obstruction classification. It follows the gravity-as-encoding
paper and precedes the quantization paper, in which nil obstructions are shown
to force discrete constraint encodings. Realization-specific constructions of
gauge bundles and connections are deferred to the geometric realization papers.
\end{remark}

% ============================================================
% Next section will be:
% \section{Lens Obstructions as Non-Unique Local Lifts}
% ============================================================
\section{Lens Obstructions as Non-Unique Local Lifts}

In this section we formalize lens obstructions in a manner suitable for defining
gauge structure as an encoding class. The key idea is that lens obstructions
correspond to non-uniqueness of admissible local lifts of coherent structure,
without contradiction or failure of representability.

\subsection{Local lifts and redundancy}

Let $\mathcal E_\alpha = (A_\alpha, Z_\alpha, E_\alpha, \ldots)$ be an admissible
encoding on domain $A_\alpha$.

\begin{definition}[Local lift]
A \emph{local lift} of a coherent structure on $A_\alpha$ is a choice of
representative in the encoding fiber $Z_\alpha$ consistent with admissibility
and overlap constraints.
\end{definition}

Local describability requires the existence of at least one admissible local
lift on each admissible domain. Lens obstructions arise when such lifts are not
unique.

\subsection{Definition of lens obstruction (refined)}

We now restate lens obstructions in terms appropriate for gauge structure.

\begin{definition}[Lens obstruction (structural)]
A \emph{lens obstruction} exists on an admissible domain if the encoding fibers
admit a nontrivial automorphism (isotropy) group acting on local lifts such that:
\begin{enumerate}[label=(\roman*)]
\item multiple admissible local lifts of the same coherent structure exist;
\item these lifts are related by admissible fiber automorphisms;
\item the automorphism structure persists under admissible refinement and cannot
be absorbed into overlap re-encoding.
\end{enumerate}
\end{definition}


Thus, lens obstructions represent redundancy of representation without
inconsistency.
\begin{remark}[Lens versus gauge]
A lens obstruction is a \emph{structural condition}: the existence of non-unique
but admissible local lifts related by persistent isotropy of encoding fibers.
Gauge structure is not identified with the lens obstruction itself, but with the
\emph{encoding response} required to organize and compensate for that
redundancy. Thus, lens denotes the obstruction, while gauge denotes the encoding
that resolves it. This distinction is structural and holds independently of any
particular physical realization.
\end{remark}

\subsection{Lens versus circle and nil}

It is important to distinguish lens obstructions sharply from the other two
obstruction types.

\begin{itemize}
\item Lens obstructions do not involve path dependence; transporting different
lifts along admissible continuation chains yields equivalent results.
\item Lens obstructions do not involve termination; admissible descriptions
exist everywhere in the domain.
\end{itemize}

\begin{remark}
If non-uniqueness of lifts could be eliminated by refining the admissible cover
or by admissible re-encoding, then the obstruction would not be a genuine lens.
Lens obstructions are precisely those redundancies that persist under all such
operations.
\end{remark}

\subsection{Fiberwise characterization}

Lens obstructions are fiberwise in nature.

\begin{lemma}
Lens obstructions correspond to nontrivial automorphism structure of encoding
fibers.
\end{lemma}

\begin{proof}
Multiple admissible local lifts correspond to distinct points in the same
encoding fiber that are related by admissible re-encoding and produce identical
coherent content. The set of such transformations forms a nontrivial automorphism
group acting on the fiber.
\end{proof}

This automorphism structure is the structural origin of gauge freedom.

\subsection{Absence of global lift}

Lens obstructions prevent the existence of a global lift.

\begin{theorem}[Non-existence of a global lift]
If a lens obstruction exists on an admissible domain, then no globally defined
admissible lift exists that is compatible with all local encodings.
\end{theorem}

\begin{proof}
A global lift would select a single representative in each encoding fiber.
However, the existence of a persistent nontrivial automorphism group implies that
any such choice can be transformed into an inequivalent admissible lift.
Therefore no choice is invariant under admissible re-encoding, and no global lift
can exist.
\end{proof}


\subsection{Interpretation}

We emphasize that lens obstructions do not indicate ambiguity or incompleteness
of the underlying system.

\begin{remark}
Lens obstructions encode descriptive redundancy, not physical indeterminacy.
They indicate that multiple representations are equally valid descriptions of
the same coherent structure.
\end{remark}

This redundancy must be organized to preserve overlap consistency and kinematic
persistence. In the next section we show that this organization uniquely defines
gauge transformations as admissible automorphisms of encoding fibers.

% ============================================================
% Next section will be:
% \section{Gauge Transformations as Fiber Automorphisms}
% ============================================================
\section{Gauge Transformations as Fiber Automorphisms}

We now formalize gauge transformations in the Modal Triplet Theory framework.
Gauge transformations are not postulated as symmetries of an underlying physical
space; they arise as the natural automorphisms associated with lens obstructions,
i.e.\ with non-unique but admissible local lifts of coherent structure.

\subsection{Automorphisms of encoding fibers}

Let $\mathcal E_\alpha = (A_\alpha, Z_\alpha, E_\alpha, \ldots)$ be an admissible
encoding, and let $z \in Z_\alpha$ be a local lift of a coherent structure.

\begin{definition}[Fiber automorphism]
A \emph{fiber automorphism} at $\mathcal E_\alpha$ is an admissible map
\[
g_\alpha : Z_\alpha \to Z_\alpha
\]
such that:
\begin{enumerate}[label=(\roman*)]
\item $g_\alpha$ preserves coherent content, i.e.\ $E_\alpha^{-1}(z)$ and
$E_\alpha^{-1}(g_\alpha(z))$ represent the same coherent structure;
\item $g_\alpha$ is compatible with admissible re-encoding on overlaps;
\item $g_\alpha$ is invertible up to admissible equivalence.
\end{enumerate}
\end{definition}

The set of all such automorphisms forms a group under composition.

\subsection{Gauge transformations}

We now identify gauge transformations with fiber automorphisms.

\begin{definition}[Gauge transformation]
A \emph{gauge transformation} is a fiber automorphism arising from a lens
obstruction, i.e.\ an admissible transformation that maps one local lift of a
coherent structure to another equally admissible lift without altering any
observable or kinematic content.
\end{definition}

Gauge transformations therefore act on representations, not on the underlying
coherent structures themselves.
\subsection{Gauge encoding versus physical realizations}

Gauge structure, as defined in this paper, is an encoding class that resolves
lens obstructions by organizing redundancy of representation. It does not
correspond to any particular physical interaction.

\begin{remark}
Electromagnetism, weak interactions, and strong interactions are \emph{distinct
realizations} of gauge encoding, characterized by different gauge groups,
representations, and dynamics. The existence of gauge encoding does not imply the
existence of any particular force; it implies only that some redundancy
bookkeeping must exist wherever lens obstructions are present.
\end{remark}

\begin{remark}
Different coherent structures may admit different local isotropy groups, and
there is no requirement that a single gauge group act universally on all
structures. Universality of gauge encoding refers to the inevitability of
redundancy bookkeeping when lens obstructions exist, not to universality of a
specific gauge interaction.
\end{remark}

\subsection{Local nature of gauge freedom}

Gauge transformations are defined locally on admissible domains.

\begin{remark}
Gauge freedom is local because lens obstructions are local: redundancy of lifts
is detected within encoding fibers on admissible domains. Global gauge
transformations exist only when compatible local automorphisms can be chosen on
all domains, which is generically obstructed.
\end{remark}

This locality is structural, not imposed.

\subsection{Gauge equivalence}

Gauge transformations induce an equivalence relation on local lifts.

\begin{definition}[Gauge equivalence]
Two local lifts $z_1, z_2 \in Z_\alpha$ are \emph{gauge equivalent} if there exists
a gauge transformation $g_\alpha$ such that $z_2 = g_\alpha(z_1)$.
\end{definition}

Physical (coherent) content is invariant under gauge equivalence.

\subsection{Gauge invariants}

Because gauge transformations encode redundancy, only gauge-invariant quantities
are meaningful descriptors of coherent structure.

\begin{definition}[Gauge-invariant quantity]
A \emph{gauge-invariant quantity} is a function or diagnostic on encoding fibers
that is constant on gauge-equivalence classes.
\end{definition}

\begin{remark}
Gauge invariants arise naturally as those features of the encoding that descend
to the quotient by fiber automorphisms. This quotient is the maximal reduced
description compatible with lens obstructions.
\end{remark}

\subsection{Distinction from gravity}

We emphasize the distinction between gauge and gravity encodings.

\begin{remark}
Gauge transformations resolve redundancy of representation (lens obstructions).
They do not resolve path-dependent inconsistency of continuation (circle
obstructions). Conversely, the kinematic consistency encoding identified as
gravity resolves circle obstructions but does not eliminate redundancy of local
lifts. The two encodings address distinct structural problems and must not be
conflated.
\end{remark}

\subsection{Uniqueness of gauge encoding}

We now show that gauge structure is the unique encoding response to lens
obstructions.
\begin{theorem}[Universality of redundancy encoding]
Any admissible resolution of lens obstructions must factor through the quotient
groupoid obtained by modding encoding fibers by their isotropy (automorphism)
action. Consequently, any such resolution is equivalent, up to admissible
re-encoding, to a redundancy encoding implemented by fiber automorphisms.
\end{theorem}
\begin{proof}
Any admissible redundancy-resolution functor must identify all gauge-equivalent
lifts while preserving overlap consistency. This is precisely the universal
property of the quotient groupoid defined by fiber isotropy. Therefore any such
resolution factors uniquely through this quotient and is equivalent to the gauge
(redundancy) encoding.
\end{proof}
\begin{remark}[Universality of gauge revisited]
Universality of gauge structure in this context means that whenever lens
obstructions exist, some redundancy encoding must be present locally. It does
not mean that a single gauge group acts on all coherent structures. Different
lens obstructions may induce different local isotropy groups in different
encodings.
\end{remark}

\begin{theorem}[Uniqueness of gauge encoding]
Any admissible encoding that resolves lens obstructions while preserving overlap
consistency and kinematic persistence is equivalent, up to admissible
re-encoding, to a redundancy encoding implemented by fiber automorphisms.
\end{theorem}

\begin{proof}
Let $\mathcal G$ be an encoding resolving lens obstructions. By definition, it
must identify all admissible local lifts as equivalent representations of the
same coherent structure. This identification defines an equivalence relation on
encoding fibers generated by admissible automorphisms. Any alternative encoding
that preserves coherent content and overlap consistency must factor through this
equivalence. Therefore $\mathcal G$ is equivalent, up to admissible re-encoding,
to a redundancy encoding realized by fiber automorphisms.
\end{proof}

\subsection{Preview: gauge fixing and connections}

While gauge transformations encode redundancy, one may choose specific
representatives for calculational or practical purposes.

\begin{remark}
Gauge fixing corresponds to a non-canonical choice of local section within a
gauge-equivalence class. Such choices do not alter the underlying encoding
structure and may fail globally due to lens obstructions.
\end{remark}

In the next section we analyze gauge fixing and show how, in smooth realization
classes, gauge redundancy is organized by bundle structure and connection-like
objects distinct from gravitational connections.

% ============================================================
% Next section will be:
% \section{Gauge Fixing and Redundancy Bookkeeping}
% ============================================================
\section{Gauge Fixing and Redundancy Bookkeeping}

In this section we analyze gauge fixing within the redundancy encoding framework
and show how bookkeeping of redundancy gives rise to connection-like structures
in smooth realization classes. We emphasize throughout that these structures are
not gravitational: they resolve redundancy (lens obstructions), not kinematic
path dependence (circle obstructions).

\subsection{Gauge fixing as section selection}

Gauge fixing corresponds to selecting a representative from each
gauge-equivalence class.

\begin{definition}[Gauge fixing]
A \emph{gauge fixing} on an admissible domain $A_\alpha$ is a choice of local
section
\[
s_\alpha : A_\alpha \to Z_\alpha
\]
such that $s_\alpha(x)$ selects a single representative from each
gauge-equivalence class in the encoding fiber over $x$.
\end{definition}

Gauge fixing is therefore a choice of description, not a structural operation.

\subsection{Non-canonicity of gauge fixing}

We now show that gauge fixing is generically non-canonical.

\begin{theorem}[Non-existence of global gauge fixing]
If a lens obstruction exists, then no globally admissible gauge fixing exists.
\end{theorem}

\begin{proof}
A global gauge fixing is a global section selecting one representative per
gauge-equivalence class. Such a section would be invariant under all fiber
automorphisms. This contradicts the existence of a nontrivial isotropy group
acting on the fibers. Hence no global gauge fixing exists.
\end{proof}


\begin{remark}
Local gauge fixing may exist on individual admissible domains, but cannot be
extended consistently across the entire atlas.
\end{remark}

\subsection{Gauge transitions on overlaps}

On overlaps of admissible domains, different gauge fixings are related by gauge
transformations.

\begin{definition}[Gauge transition]
Let $s_\alpha$ and $s_\beta$ be gauge fixings on overlapping admissible domains
$A_\alpha$ and $A_\beta$. The \emph{gauge transition} on $A_{\alpha\beta}$ is the
gauge transformation $g_{\alpha\beta}$ such that
\[
s_\beta(x) = g_{\alpha\beta}(x)\, s_\alpha(x)
\]
for all $x \in A_{\alpha\beta}$.
\end{definition}

Gauge transitions encode how different local gauge choices are related.

\subsection{Redundancy bookkeeping}

Because gauge fixing is non-canonical, redundancy must be tracked explicitly.

\begin{remark}
Redundancy bookkeeping assigns data to overlaps that records how local gauge
choices differ. This bookkeeping ensures consistency of reduced descriptions
without privileging any particular gauge.
\end{remark}

This bookkeeping role is the defining function of gauge structure in MTT.

\subsection{Smooth realizations and gauge connections}

We now show how redundancy bookkeeping is realized geometrically in smooth
settings.

\begin{theorem}[Gauge connection realization]
In smooth realization classes, redundancy bookkeeping is realized by a
connection-like object on a principal bundle whose structure group is the gauge
automorphism group of encoding fibers.
\end{theorem}

\begin{proof}
Gauge transitions $g_{\alpha\beta}(x)$ define smooth maps on overlaps. Consistency
of these transitions under refinement and composition requires infinitesimal
control of how gauge choices vary. The unique structure implementing such control
is a connection on the associated principal bundle. This connection records how
to compare gauge choices infinitesimally without selecting a preferred global
representative.
\end{proof}
\begin{remark}[Electromagnetism as a mixed realization]
Electromagnetism provides an important example in which lens and circle
obstructions coexist within a single realization. The U(1) phase redundancy of
electromagnetism realizes a lens obstruction resolved by gauge encoding, while
the associated holonomy around closed loops realizes a circle obstruction. The
present separation of gauge and gravity reflects a distinction of structural
roles, not a claim that physical theories realize only one obstruction type at a
time.
\end{remark}

\subsection{Distinction from gravitational connection}

We stress the distinction between gauge connections and gravitational
connections.

\begin{remark}
Gauge connections arise to track redundancy of representation (lens obstructions).
Gravitational connections arise to resolve kinematic path dependence (circle
obstructions). Although both are realized geometrically as connections, they
encode distinct structural roles and should not be conflated.
\end{remark}

This distinction is structural and persists independently of particular
realization choices.

\subsection{Gauge invariance revisited}

Gauge invariance is now seen as invariance under redundancy bookkeeping.

\begin{remark}
Gauge invariance expresses the fact that physical (coherent) content is invariant
under changes of representative within a gauge-equivalence class. It is not a
symmetry of the underlying system, but a statement about redundancy of
description.
\end{remark}

\subsection{Preview: coupling to other encodings}

Gauge structure may interact with other encoding classes in particular
realizations.

\begin{remark}
In unified or saturated encodings, gauge redundancy bookkeeping may couple to
kinematic consistency bookkeeping (gravity) or to discrete constraint encodings
(quantization). Such couplings are realization-dependent and do not alter the
structural separation of encoding roles established here.
\end{remark}

\begin{remark}[Circle realizations and gravity]
Circle obstructions denote loop-dependent failure of global coherence. Different
encodings may realize this obstruction in different regimes. When circle affects
internal or phase transport, electromagnetism provides a minimal realization.
When circle affects kinematic persistence itself—worldlines, causal cones, and
horizons—a gravity-like kinematic consistency encoding is forced. These are
distinct realizations of the same structural invariant.
\end{remark}

% ============================================================
% Next section will be:
% \section{Summary and Outlook}
% ============================================================
\section{Summary and Outlook}

In this paper we have identified gauge structure as a necessary encoding class
arising from lens obstructions in the Modal Triplet Theory framework. Lens
obstructions correspond to non-unique but admissible local representations of
coherent structure, with full overlap consistency and kinematic persistence.
Such redundancy cannot be eliminated by refinement or admissible re-encoding and
therefore requires explicit bookkeeping.

We showed that this bookkeeping is uniquely implemented by a redundancy encoding
whose transformations act as automorphisms of encoding fibers. Gauge
transformations are thus not symmetries of an underlying physical arena, but
automorphisms of representation that relate equally valid local lifts. Gauge
equivalence expresses invariance of coherent content under changes of
representative, and gauge fixing corresponds to a non-canonical choice of local
section that cannot be extended globally in the presence of lens obstructions.

In smooth realization classes, redundancy bookkeeping is realized geometrically
by principal bundles equipped with connection-like objects that track transitions
between local gauge choices. These gauge connections are structurally distinct
from the gravitational connections introduced to resolve circle obstructions.
Although both appear as connections in geometric realizations, they encode
different kinds of bookkeeping: gauge resolves redundancy of representation,
while gravity resolves path-dependent kinematic inconsistency.

The results of this paper complete the identification of gauge structure within
the Modal Triplet Theory Program. Together with the gravity-as-encoding paper,
they establish that two of the three fundamental obstruction types—circle and
lens—require distinct and unavoidable encoding responses. The remaining
obstruction type, nil, corresponds to termination of describability and motivates
a discrete constraint encoding, treated in the subsequent paper on quantization.

More broadly, this analysis clarifies why gauge structure is universal yet
unobservable directly, why gauge fixing is inherently non-canonical, and why
gauge symmetry appears as a freedom of description rather than as a property of
underlying reality. These features are not postulated; they are forced by the
structure of local describability and overlap consistency.

Subsequent papers in the series develop the encoding response to nil obstructions
(quantization), analyze the intersection and saturation of multiple encoding
classes (including the Standard Model and string-theoretic frameworks), and
construct explicit realizations in geometric and bundle-based models. No
additional structural assumptions are introduced downstream.

In this way, gauge structure emerges as an inevitable component of any coherent
descriptive framework once redundancy of representation is present, completing
the lens branch of the circle–lens–nil triad.

\begin{remark}[Universality clarified]
Universality of gauge structure in the Modal Triplet Theory framework means that
whenever lens obstructions exist, some redundancy encoding must exist locally.
It does not mean that a single gauge group, interaction, or coupling acts on all
coherent structures. Different realizations may exhibit different gauge groups,
or none at all, depending on the presence and nature of lens obstructions.
\end{remark}


% ============================================================
% End of Paper B4

\end{document}
